\documentclass[aspectratio=169,11pt]{beamer}

\usepackage[utf8]{inputenc}
\usepackage[T1]{fontenc}
\usepackage[spanish]{babel}
\usepackage{graphicx}
\usepackage{booktabs}
\usepackage{amsmath}
\usepackage{hyperref}

\usetheme{Madrid}
\usecolortheme{default}
\setbeamertemplate{navigation symbols}{}
\setbeamertemplate{caption}[numbered]

\title[ACP calidad del agua]{Análisis de Componentes Principales aplicado a calidad del agua del Río Cauca}
\subtitle{Taller 1. Análisis Multivariado}
\author[J. Tapia, J. Jaramillo, L. Meza]{
JHONATAN ANDRES TAPIA CORDOBA\\
JORGE ANDRES JARAMILLO NEME\\
LUIS FERNANDO MEZA RAMREZ
}
\institute{Maestría en Inteligencia Artificial y Ciencia de Datos}
\date{Mayo 2 de 2026}

\newcommand{\img}[2]{
  \begin{center}
    \includegraphics[width=#2\textwidth,height=0.62\textheight,keepaspectratio]{\detokenize{#1}}
  \end{center}
}

\begin{document}

\begin{frame}
  \titlepage
\end{frame}

\begin{frame}{Objetivo de la sustentación}
  \begin{itemize}
    \item Aplicar ACP sobre variables físico-químicas de calidad del agua del Río Cauca.
    \item Resumir la información multivariada en dimensiones interpretables.
    \item Interpretar los principales gradientes ambientales.
    \item Identificar observaciones atípicas.
    \item Evaluar si es posible construir un índice.
  \end{itemize}

  \vspace{0.3cm}
  \textbf{Importante:} el ACP permite pasar de varias variables correlacionadas a pocas dimensiones interpretables.
\end{frame}

\begin{frame}{Datos y preparación}
  \begin{itemize}
    \item Base original: \textbf{2368 registros y 56 variables}.
    \item Se limpiaron nombres de columnas y formatos numéricos.
    \item Se corrigieron valores con coma decimal y notación científica tipo \texttt{*10E}.
    \item Se construyó el identificador \texttt{anio\_estacion}.
    \item Las variables con 100 \% de valores faltantes fueron descartadas.
  \end{itemize}

  \vspace{0.25cm}
  Cada individuo se interpreta como una observación espacio-temporal: año más estación de monitoreo.
\end{frame}

\begin{frame}{Variables activas y tratamiento}
  \begin{columns}[T]
    \begin{column}{0.48\textwidth}
      \begin{itemize}
        \item Turbiedad.
        \item Sólidos suspendidos totales.
        \item DBO.
        \item DQO.
        \item Oxígeno disuelto.
        \item Conductividad eléctrica.
        \item Fósforo total.
      \end{itemize}
    \end{column}
    \begin{column}{0.48\textwidth}
      \begin{itemize}
        \item Todas presentaron asimetría positiva severa.
        \item Se aplicó transformación \texttt{log1p}.
        \item Luego se estandarizaron por estar en unidades diferentes.
        \item Matriz final del ACP: \textbf{1475 observaciones}.
        \item Conservación: \textbf{62,29 \%}.
      \end{itemize}
    \end{column}
  \end{columns}
\end{frame}

\begin{frame}{Estadísticas descriptivas}
  \begin{itemize}
    \item La exploración descriptiva mostró alta heterogeneidad.
    \item Los coeficientes de variación fueron altos, especialmente en DBO, fósforo total y oxígeno disuelto.
    \item Esto justificó la transformación logarítmica antes del ACP.
  \end{itemize}

  \img{Graficos de resultados/distribuciones fisico-quimica_grupo_A.png}{0.95}
\end{frame}

\begin{frame}{Distribuciones físico-químicas}
  \img{Graficos de resultados/distribuciones fisico-quimica_grupo_B.png}{0.95}

  \vspace{-0.2cm}
  \small Las distribuciones muestran asimetrías positivas y valores extremos, por lo que se usó \texttt{log1p}.
\end{frame}

\begin{frame}{Matriz de correlación}
  \begin{itemize}
    \item Turbiedad y sólidos suspendidos totales: \textbf{r = 0,80}.
    \item DBO y fósforo total: \textbf{r = 0,60}.
    \item Oxígeno disuelto y conductividad eléctrica: \textbf{r = -0,54}.
  \end{itemize}

  \img{Graficos de resultados/Matriz de correlacion: Espacion de variable.png}{0.88}
\end{frame}

\begin{frame}{Varianza explicada por el ACP}
  \begin{itemize}
    \item Criterio de Kaiser: se retienen \textbf{tres componentes}.
    \item Dim 1: \textbf{31,57 \%}.
    \item Dim 2: \textbf{25,43 \%}.
    \item Dim 3: \textbf{20,15 \%}.
    \item Acumulado en tres dimensiones: \textbf{77,15 \%}.
    \item Plano 1-2: \textbf{57,00 \%}.
  \end{itemize}
\end{frame}

\begin{frame}{Sedimentación de varianza}
  \img{Graficos de resultados/Sedimentacion de varianza.png}{0.95}

  \vspace{-0.2cm}
  \small El plano 1-2 es útil para la lectura gráfica, aunque la tercera dimensión conserva información adicional.
\end{frame}

\begin{frame}{Nube de individuos}
  \begin{itemize}
    \item Los puntos cercanos al origen tienen perfiles más promedio.
    \item Los puntos alejados representan perfiles físico-químicos diferenciados.
    \item Los individuos mejor representados tienen mayor coseno cuadrado en el plano.
  \end{itemize}

  \img{Graficos de resultados/Nube de individuos.png}{0.9}
\end{frame}

\begin{frame}{Individuos mejor representados}
  \begin{table}
    \centering
    \small
    \begin{tabular}{lrrr}
      \toprule
      Individuo & Dim 1 & Dim 2 & Cos2 plano \\
      \midrule
      2013\_ANTES SUAREZ & -2,145 & 1,375 & 0,992 \\
      2011\_PASO DEL COMERCIO & 1,929 & 1,695 & 0,985 \\
      1997\_PUENTE HORMIGUERO & -2,190 & 1,341 & 0,978 \\
      1998\_ANTES RIO OVEJAS & -4,093 & -0,461 & 0,976 \\
      2000\_PASO DE LA BOLSA & -1,714 & 1,469 & 0,972 \\
      \bottomrule
    \end{tabular}
  \end{table}

  \small Estos casos tienen representación confiable en el plano factorial 1-2.
\end{frame}

\begin{frame}{Círculo de correlaciones}
  \begin{itemize}
    \item \textbf{Dim 1:} sólidos suspendidos, turbiedad, DBO, DQO y fósforo total.
    \item Interpretación: gradiente de carga contaminante, orgánica y particulada.
    \item \textbf{Dim 2:} oxígeno disuelto y conductividad eléctrica en sentidos opuestos.
    \item Interpretación: contraste entre oxigenación e influencia iónica/conductiva.
  \end{itemize}

  \img{Graficos de resultados/Circulos de correlaciones Plano 1-2.png}{0.9}
\end{frame}

\begin{frame}{Biplot y datos atípicos}
  \begin{itemize}
    \item El biplot integra individuos y variables.
    \item Las observaciones en dirección de una variable tienden a presentar valores altos en ella.
    \item Los posibles atípicos se identificaron por distancia al origen en el plano factorial.
  \end{itemize}

  \img{Graficos de resultados/Relacion Dual individuos - variables.png}{0.88}
\end{frame}

\begin{frame}{Posibles atípicos factoriales}
  \begin{columns}[T]
    \begin{column}{0.48\textwidth}
      \begin{itemize}
        \item 1993\_JUANCHITO.
        \item 2010\_MEDIACANOA.
        \item 2010\_PASO DEL COMERCIO.
        \item 2010\_PUERTO ISAACS.
        \item 2010\_JUANCHITO.
        \item 2010\_YOTOCO.
      \end{itemize}
    \end{column}
    \begin{column}{0.48\textwidth}
      \small
      Estos registros no deben eliminarse automáticamente. Pueden representar eventos reales de contaminación, cambios hidrológicos o condiciones ambientales particulares.
    \end{column}
  \end{columns}
\end{frame}

\begin{frame}{Síntesis con contribuciones y cosenos}
  \begin{itemize}
    \item Aporte promedio esperado por variable: \textbf{14,29 \%}.
    \item En Dim 1 dominan:
      \begin{itemize}
        \item Sólidos suspendidos totales: \textbf{32,36 \%}.
        \item Turbiedad: \textbf{23,45 \%}.
        \item DBO: \textbf{15,71 \%}.
      \end{itemize}
    \item En Dim 2 dominan:
      \begin{itemize}
        \item Oxígeno disuelto: \textbf{41,02 \%}.
        \item Conductividad eléctrica: \textbf{32,98 \%}.
      \end{itemize}
  \end{itemize}
\end{frame}

\begin{frame}{Construcción del índice}
  \begin{itemize}
    \item Sí es posible construir un índice, pero como \textbf{índice parcial}.
    \item No debe interpretarse como índice global de calidad del agua.
    \item PC1 explica \textbf{31,57 \%} y resume carga contaminante/particulada.
    \item PC2 conserva información importante sobre oxígeno disuelto y conductividad.
  \end{itemize}

  \[
    \text{Indice\_PC1\_0\_100} =
    100 \times \frac{PC1 - \min(PC1)}{\max(PC1) - \min(PC1)}
  \]
\end{frame}

\begin{frame}{Resultados del índice parcial}
  \begin{itemize}
    \item Valores cercanos a 100 indican mayor carga contaminante/particulada según PC1.
    \item Mayores valores individuales:
      \begin{itemize}
        \item 2010\_LA VICTORIA.
        \item 1994\_PASO DEL COMERCIO.
        \item 2010\_ANACARO.
        \item 2010\_MEDIACANOA.
        \item 2010\_PASO DEL COMERCIO.
      \end{itemize}
    \item Mayores promedios por estación: ANACARO, PUENTE LA VIRGINIA, LA VICTORIA, VIJES y PUENTE GUAYABAL.
  \end{itemize}
\end{frame}

\begin{frame}{Conclusiones}
  \begin{enumerate}
    \item El ACP resumió siete variables físico-químicas en dimensiones interpretables.
    \item La Dim 1 representa carga contaminante y particulada.
    \item La Dim 2 complementa la lectura mediante el contraste entre oxígeno disuelto y conductividad eléctrica.
    \item Se identificaron observaciones extremas potencialmente relevantes.
    \item El índice construido debe leerse como índice parcial de carga contaminante.
  \end{enumerate}

  \vspace{0.3cm}
  \textbf La calidad del agua del Río Cauca no se resume completamente en un solo eje.
\end{frame}

\begin{frame}{Preguntas probables}
  \small
  \textbf{¿Por qué log1p?} Por asimetría positiva severa y valores extremos.\\[0.25cm]
  \textbf{¿Por qué estandarizar?} Porque las variables estaban en unidades distintas.\\[0.25cm]
  \textbf{¿Por qué tres componentes?} Porque tres valores propios fueron mayores que 1 y explicaron 77,15 \%.\\[0.25cm]
  \textbf{¿Por qué el índice no es global?} Porque PC2 contiene información ambiental relevante no resumida por PC1.\\[0.25cm]
  \textbf{¿Qué significa un atípico en ACP?} Una observación alejada del centro factorial; no necesariamente un error.
\end{frame}

\end{document}
